\section{Visual Substitution Alphabets}\label{sec:cipher}

\subsection{Definitions and a worked example}

Let $\Sigma=\{\mathrm{A},\dots,\mathrm{Z}\}$ be the Latin capitals and $\Sigma^{*}$ the set of finite strings over $\Sigma$. A \emph{substitution} is a bijection $\pi:\Sigma\to\Sigma$, extended letter by letter to a map on strings, $\pi(x_1x_2\cdots x_n)=\pi(x_1)\pi(x_2)\cdots\pi(x_n)$. Characters outside $\Sigma$, namely spaces and punctuation, are left fixed, and case is preserved by applying $\pi$ to the upper-case counterpart of a lower-case letter and returning the result in lower case. The substitution is \emph{reciprocal} when $\pi\circ\pi=\id$, which is the same as saying that $\pi$ is an involution: it exchanges letters in pairs and leaves the others where they are. A reciprocal substitution is its own decoder, and applying it twice restores the original text.

A \emph{visual} substitution is one chosen so that $\pi(a)$ resembles $a$ for most letters $a$. The notion of resemblance is left open for the moment and is taken up in Section~\ref{sec:resemblance}. Consider the exchanges B$\leftrightarrow$P and U$\leftrightarrow$V and nothing else. The word PUBLIC becomes BVPLIC: the P and B exchange, the U becomes V, and the remaining letters are untouched. BLUE becomes PLVE and LOVE becomes LOUE, and in each case a second application restores the original, so that PUBLIC passes through BVPLIC and returns. These are small examples, but they are enough to show the features that matter. The result is close to the original in shape, a sentence supplies cues that disambiguate an individual replacement, and a reader who meets the substituted spelling a few times may come to experience it as an unfamiliar typeface or an alternative orthography rather than as a code to be solved. The same examples show a hazard of hand application: it is easy to leave the U of PUBLIC in place while exchanging the letters around it, and a substitution of even six letters is error-prone when applied by eye. This is a modest argument for generating substituted text mechanically from a stored source, a point that returns in Section~\ref{sec:edition}.

\subsection{Counting}

The number of substitutions of the 26 capitals is $26!\approx4.03\times10^{26}$. The number of reciprocal substitutions is far smaller. If $I(n)$ denotes the number of involutions on an $n$-element set, then letter $n$ either is fixed, leaving $I(n-1)$ involutions of the rest, or is exchanged with one of the other $n-1$ letters, leaving $I(n-2)$ involutions of the remaining $n-2$, so that
\[
I(n)=I(n-1)+(n-1)\,I(n-2),\qquad I(0)=I(1)=1.
\]
A reciprocal substitution that changes exactly $2k$ letters is determined by choosing which $2k$ letters move and then a perfect matching on them, and the number of perfect matchings on $2k$ points is $(2k-1)!!=(2k-1)(2k-3)\cdots 3\cdot 1$. Hence the number of reciprocal substitutions that exchange exactly $k$ pairs is $\binom{26}{2k}(2k-1)!!$, and summing over $k$ gives $I(26)$. Running these formulas gives the figures in Table~\ref{tab:counting}; the sum over $k$ reproduces $I(26)=532{,}985{,}208{,}200{,}576$ exactly, which serves as a check on the arithmetic. The count of reciprocal substitutions that change every letter, $25!!=7{,}905{,}853{,}580{,}625$, is the number of perfect matchings on the full alphabet, and the count of all substitutions that change every letter is the derangement number $D(26)=148{,}362{,}637{,}348{,}470{,}135{,}821{,}287{,}825$, which satisfies $D(n)=(n-1)(D(n-1)+D(n-2))$.

\begin{table}[ht]
\centering
\small
\begin{tabular}{rrr}
\toprule
Pairs exchanged $k$ & Letters changed $2k$ & Reciprocal alphabets $\binom{26}{2k}(2k-1)!!$\\
\midrule
1 & 2 & 325\\
2 & 4 & 44\,850\\
3 & 6 & 3\,453\,450\\
5 & 10 & 5\,019\,589\,575\\
8 & 16 & 10\,767\,019\,638\,375\\
10 & 20 & 150\,738\,274\,937\,250\\
13 & 26 & 7\,905\,853\,580\,625\\
\midrule
\multicolumn{2}{l}{All reciprocal alphabets, $I(26)$} & 532\,985\,208\,200\,576\\
\multicolumn{2}{l}{All permutations without fixed points, $D(26)$} & 148\,362\,637\,348\,470\,135\,821\,287\,825\\
\multicolumn{2}{l}{All permutations, $26!$} & 403\,291\,461\,126\,605\,635\,584\,000\,000\\
\bottomrule
\end{tabular}

\caption{Reciprocal substitutions of the 26 Latin capitals by the number of exchanged pairs, with the totals for all reciprocal substitutions, all substitutions without fixed points, and all substitutions.}\label{tab:counting}
\end{table}

The ratio of all permutations to reciprocal ones is about $7.6\times10^{11}$, so restricting to involutions discards all but a vanishing fraction of the space. Two further observations are worth recording. First, the space is enormous in the sense relevant to cryptography and irrelevant to the present purpose, since a fixed substitution alphabet falls to frequency analysis and word-pattern analysis however the shapes are paired. The security parameter of a visual substitution is deliberately negligible, because the aim is that the reader's eye should be able to see through it. Second, the space that matters is a thin slice of the whole. A substitution that makes every letter resemble a different letter, with a resemblance high enough to be useful, is constrained by the geometry of the shapes, and the effective design space is the set of alphabets that satisfy those constraints, a small subset of the figures in the table. The interesting design problem is therefore not the size of the space but the identification of its good regions.

\subsection{Conventions that support reading}

A visual substitution is not meant to be read letter by letter. Skilled reading draws on word shape, on the redundancy of the language, and on the predictability of words from their context. Shannon's experiments on the prediction of printed English put the information content of ordinary text at roughly one bit per letter, against the $\log_2 26\approx4.7$ bits that a uniformly random letter would carry \citep{shannon1951}. A reader therefore has on average a great deal of contextual help with any one letter, and a substitution that preserves spaces, punctuation and case leaves that help intact. This is the sense in which the ambiguity of an individual replacement is tolerable: the replacement of a letter by a lookalike changes a small fraction of a word's features, and the lexicon plus the sentence usually settles the rest. It is also why the quality of an alphabet should be judged on whole passages. Several individually plausible exchanges can combine into something confusing, for instance when the same word contains two exchanged letters whose joint effect is a different familiar word, and a reciprocal pairing is best treated as a first approximation to be adjusted by trying it on ordinary prose.

\subsection{Cycles longer than two}

Reciprocal exchange is not the only structure available. A substitution may contain longer cycles, such as I$\to$L$\to$T$\to$I, in which each letter is replaced by a different neighbour and the inverse cycle is needed for decoding. The loss relative to the reciprocal case is that encoding and decoding use different tables, and the gain is a larger set of shape-compatible alphabets, because a cycle needs only that each letter resemble its successor and not that the relation hold in both directions. A further and less obvious difference between cycles of length two and cycles of length three or more appears when the shapes are interpolated, and it is established in Section~\ref{sec:drift}: a transposition cannot be effected by linear interpolation without the two letters sharing a glyph at some moment, whereas a three-cycle in general position can.
